\documentclass[crop,tikz,11pt]{standalone}
\usepackage{geometricfigures}
\usepackage{amsmath,amssymb}
\usetikzlibrary{arrows.meta}

\definecolor{lib}{HTML}{3B75AF}
\definecolor{rot}{HTML}{C1701E}
\definecolor{sep}{HTML}{B03030}
\definecolor{lobe}{HTML}{FFD700}   % the lobe K, as (:gold, 0.5) in the generated plots: fill lobe!50

% ---------------------------------------------------------------------------
% From the Morse picture to the two loops at a saddle: panels (b)-(d) next to morse-saddle,
% which is panel (a), the hyperbolic paraboloid.
%
% (b) The latent plane near a saddle P. Inside the neighborhood N_P (dashed) the level set through P
%     is the preimage under the Morse chart y of the lines y_2 = +-y_1: two smooth curves crossing
%     transversally at P. The chart is drawn as y^{-1}(y_1, y_2) = (y_1, y_2 + k y_1^2), a
%     diffeomorphism that bends the lines into curves; levels below E_P (blue) lie in the upper and
%     lower sectors, levels above (orange) in the left and right ones, as in panel (a). The four
%     half-branches are B_1, ..., B_4, counterclockwise from the upper right.
% (c), (d) Continued beyond N_P, the four half-branches at P join, adjacent ones in pairs, into two
%     loops: side by side (c), or nested (d), with the crescent U between them (blue tint) and the
%     lobe K inside the inner loop (yellow). The same four half-branches leave P in all three panels:
%     B_2 and B_3 (upper left, lower left) form gamma_-, B_4 and B_1 (lower right, upper right) gamma_+.
%
% Schematic: no values are measured.
\newcommand{\kk}{0.45}   % bend of the drawn chart
\newcommand{\cc}{0.12}  % level offset c of the hyperbolas in (b)
\begin{document}
\begin{tikzpicture}[>=Stealth,
  gam/.style={sep,very thick},
  tag/.style={font=\footnotesize,inner sep=1.5pt},
  panelcap/.style={font=\small,anchor=north}]

% ===========================================================================
% (b) near P
% ===========================================================================
\begin{scope}[scale=1.75]
  \begin{scope}
    \clip (0,0) circle (1);
    \foreach \s in {-1,1} {
      % levels E_P - c: y_2 = +-sqrt(y_1^2 + c), upper and lower sectors
      \draw[lib,thin,domain=-1:1,samples=40,smooth]
        plot ({\x},{\s*sqrt(\x*\x+\cc)+\kk*\x*\x});
      % levels E_P + c: y_1 = +-sqrt(y_2^2 + c), left and right sectors
      \draw[rot,thin,domain=-1:1,samples=40,smooth]
        plot ({\s*sqrt(\x*\x+\cc)},{\x+\kk*(\x*\x+\cc)});
    }
    % the level set through P: the preimages of y_2 = y_1 and y_2 = -y_1
    \draw[gam,domain=-1:1,samples=30,smooth] plot ({\x},{\x+\kk*\x*\x});
    \draw[gam,domain=-1:1,samples=30,smooth] plot ({\x},{-\x+\kk*\x*\x});
  \end{scope}
  % the four half-branches at P, counterclockwise from the upper right
  \node[tag,sep,anchor=south west] at (0.64,0.80) {$\mathcal{B}_1$};
  \node[tag,sep,anchor=south east] at (-0.64,0.80) {$\mathcal{B}_2$};
  \node[tag,sep,anchor=north east] at (-0.86,-0.52) {$\mathcal{B}_3$};
  \node[tag,sep,anchor=north west] at (0.86,-0.52) {$\mathcal{B}_4$};
  \draw[fg!55!bg,dashed] (0,0) circle (1);
  \node[tag,fg!70!bg,anchor=north] at (0,-1.0) {$N_P$};
  \fill[sep] (0,0) circle (1.0pt);
  \node[tag,sep,anchor=north] at (0,-0.07) {$P$};
\end{scope}

% ===========================================================================
% (c) side by side
% ===========================================================================
\begin{scope}[shift={(5.4,0)},scale=0.78]
  \draw[fg!40!bg,dashed] (0,0) circle (0.55);
  % gamma_-: upper-left and lower-left half-branches
  \draw[gam] (0,0) .. controls (-0.9,0.9) and (-2.4,1.5) .. (-2.4,0)
                   .. controls (-2.4,-1.5) and (-0.9,-0.9) .. (0,0);
  % gamma_+: upper-right and lower-right half-branches
  \draw[gam] (0,0) .. controls (0.9,0.9) and (2.4,1.5) .. (2.4,0)
                   .. controls (2.4,-1.5) and (0.9,-0.9) .. (0,0);
  \fill[sep] (0,0) circle (1.7pt);
  \node[tag,sep,anchor=north] at (0,-0.55) {$P$};
  \node[tag,sep,anchor=south] at (-1.7,1.15) {$\gamma_-$};
  \node[tag,sep,anchor=south] at ( 1.7,1.15) {$\gamma_+$};
\end{scope}

% ===========================================================================
% (d) nested
% ===========================================================================
\begin{scope}[shift={(10.6,0)},scale=0.78]
  % the crescent U between the loops, and the lobe K inside the inner one
  \fill[lib!14!bg,even odd rule]
    (0,0) .. controls (-0.9,0.9) and (-1.4,2.3) .. (0.9,2.3)
          .. controls (3.4,2.3) and (3.4,-2.3) .. (0.9,-2.3)
          .. controls (-1.4,-2.3) and (-0.9,-0.9) .. (0,0)
    (0,0) .. controls (0.9,0.9) and (2.1,1.1) .. (2.1,0)
          .. controls (2.1,-1.1) and (0.9,-0.9) .. (0,0);
  \fill[lobe!50!bg]
    (0,0) .. controls (0.9,0.9) and (2.1,1.1) .. (2.1,0)
          .. controls (2.1,-1.1) and (0.9,-0.9) .. (0,0);
  \draw[fg!40!bg,dashed] (0,0) circle (0.55);
  % gamma_-: upper-left and lower-left half-branches, around the right
  \draw[gam] (0,0) .. controls (-0.9,0.9) and (-1.4,2.3) .. (0.9,2.3)
                   .. controls (3.4,2.3) and (3.4,-2.3) .. (0.9,-2.3)
                   .. controls (-1.4,-2.3) and (-0.9,-0.9) .. (0,0);
  % gamma_+: upper-right and lower-right half-branches
  \draw[gam] (0,0) .. controls (0.9,0.9) and (2.1,1.1) .. (2.1,0)
                   .. controls (2.1,-1.1) and (0.9,-0.9) .. (0,0);
  \fill[sep] (0,0) circle (1.7pt);
  \node[tag,sep,anchor=east] at (-0.55,0) {$P$};
  \node[tag,sep,anchor=south] at (0.9,2.3) {$\gamma_-$};
  \node[tag,sep,anchor=west] at (2.1,0) {$\gamma_+$};
  \node[tag,lib] at (0.9,1.5) {$U$};
  \node[tag,fg!70!bg] at (1.35,0) {$K$};
\end{scope}

% ===========================================================================
\node[panelcap] at (0,-2.3) {\textbf{(b)}};
\node[panelcap] at (5.4,-2.3) {\textbf{(c)}};
\node[panelcap] at (10.6,-2.3) {\textbf{(d)}};
\end{tikzpicture}
\end{document}
